\section{Lower bounds and sharp rates}

Rate matching is proved on the explicit uniformly conditioned two-class specialization \(k=d_x=d_z=2\), \(L=2\), \(\pi_0=1/10\), and \(\sigma_0=1/10\). The construction uses two related two-class submodels. The first moves a collision witness away from a single class-average effect value at \(n^{-1/2}\) scale; the second changes ordered latent masses along a factorization-preserving path whose observed separation is proportional to the product of the effect gap and the weight displacement. The lower bounds combine the displayed local Kullback--Leibler neighborhoods with the standard two-point testing logic of \citet{LeCam1986} and \citet{PolyanskiyWu2019}.

\begin{definitionv}{def:local-kl-radius}[Local Kullback--Leibler radius \(c_{\mathrm{loc}}\)]
Fix a local Kullback--Leibler radius constant \(c_{\mathrm{loc}}\) satisfying
\[
0<c_{\mathrm{loc}}<1.
\]
\end{definitionv}

The constant \(c_{\mathrm{loc}}\) fixes the radius of the local experiments used below. Keeping this radius bounded away from one gives a common testing scale for the two-point comparisons.

\begin{definitionv}{synth_14}[Factorization-preserving two-class path law $P_{g,h}$]
For $0\le g\le 1/4$ and $|h|\le 1/100$, $P_{g,h}$ is the full-data law on $(U,T,X,Z,Y(0),Y(1),Y)$ with $U\in\{1,2\}$, $T\in\{0,1\}$, $X=(1,X_2)^\top$, $Z=(1,Z_2)^\top$, $Y=Y(T)$, latent masses $P_{g,h}(U=1)=2/5+h$ and $P_{g,h}(U=2)=3/5-h$, and latent-arm weights\[\begin{array}{c|cc} & u=1 & u=2\\ \hline P_{g,h}(U=u,T=1) & 4/25+3h/5 & 9/25-3h/5\\ P_{g,h}(U=u,T=0) & 6/25+2h/5 & 6/25-2h/5 .\end{array}\] Thus $P_{g,h}(T=1\mid U=u)=P_{g,h}(U=u,T=1)/P_{g,h}(U=u)$. Conditional on $U$, the variables $T,X_2,Y(0),Y(1)$ are Bernoulli with the displayed treatment probability, with $E_{P_{g,h}}[X_2\mid U=1]=1/5$, $E_{P_{g,h}}[X_2\mid U=2]=(12/25-h/5)/(3/5-h)$, $E_{P_{g,h}}[Y(0)\mid U=u]=1/4$, $E_{P_{g,h}}[Y(1)\mid U=1]=1/2-g/2$, and $E_{P_{g,h}}[Y(1)\mid U=2]=1/2+g/2$. Conditional on $(U,T)$, $Z_2$ is Bernoulli and independent of $(X_2,Y(0),Y(1))$, with\[E_{P_{g,h}}[Z_2\mid U=1,T=1]=\frac{225h+28}{20(15h+4)},\quad E_{P_{g,h}}[Z_2\mid U=2,T=1]=\frac34,\]\[E_{P_{g,h}}[Z_2\mid U=1,T=0]=\frac{35h+9}{10(5h+3)},\quad E_{P_{g,h}}[Z_2\mid U=2,T=0]=\frac7{10}.\] With the induced feature matrices $A_t(h)$ and $B(h)$ and latent-arm weight vector $w_t(h)=(P_{g,h}(U=1,T=t),P_{g,h}(U=2,T=t))$, the path preserves the armwise proxy factorization $A_t(h)\operatorname{diag}(w_t(h))B(h)^\top=A_t(0)\operatorname{diag}(w_t(0))B(0)^\top$ for $t\in\{0,1\}$, and $P_{g,0}=P^{\mathrm{wit}}_{g/2}$.
\end{definitionv}

\begin{definitionv}{def:path-law-family}[Path law \(P_{g,h}\)]
For a gap \(g\) and displacement \(h\), \(P_{g,h}\) is the explicit two-class factorization-preserving path law obtained by perturbing the latent weights while preserving the stated proxy factorizations.
\end{definitionv}

The divergence \(D_{\mathrm{KL}}\) is used only for the observed one-record margins that generate the sample laws \(Q_P^{(n)}\). Thus the lower bounds compare statistical experiments through the same observed records as the estimators and confidence reports.

The local neighborhoods center the two testing problems at the relevant collision or separated path endpoint.

\begin{definitionv}{def:published-vmw-regimes}[Published VMW scope handle]
The published VMW scope handle is the cited-publication interface for the proxy-model parameter class and recovery regime. For fixed latent cardinality \(k\), target-proxy dimension \(d_x\), reference-proxy dimension \(d_z\), and full-data law \(P\) on \((U,T,X,Z,Y(0),Y(1),Y)\), it records three law-level predicates:
\[
\mathsf{Model}_{\mathrm{VMW}}(k,d_x,d_z,P),
\qquad
\mathsf{FiniteReg}_{\mathrm{VMW}}(k,d_x,d_z,P),
\qquad
\mathsf{Recover}_{\mathrm{VMW}}(k,d_x,d_z,P).
\]
These predicates are publication-level gates whose concrete interpretations enter through separate cited gates. The handle also includes the published estimator
\[
\widehat\Theta^{\mathrm{VMW}}_n:
(O_{d_x,d_z})^n
\to
\mathbb R^k\times\mathbb R^{d_x\times k}\times\mathbb R^k,
\]
which maps an observed sample of size \(n\) from records \((T,X,Z,Y)\) to a triple consisting of treatment effects, an anchor-normalized target-proxy feature matrix, and real-valued weight coordinates.
\end{definitionv}

The quotient-law lower bound applies the Le Cam local testing framework of \citet{LeCam1986} and \citet{PolyanskiyWu2019} to \cref{obj:ass:local-quotient-neighborhood}. It keeps the observed experiment within \(n^{-1}\) divergence of the collision witness, so an \(n\)-sample comparison has bounded information distance.

\begin{definitionv}{synth_6}[Kullback--Leibler divergence $D_{\mathrm{KL}}$]
For probability laws $\mu$ and $\nu$ on the same measurable space, define $D_{\mathrm{KL}}(\mu\|\nu):=\int \log(d\mu/d\nu)\,d\mu$ when $\mu\ll\nu$, and set $D_{\mathrm{KL}}(\mu\|\nu):=+\infty$ when $\mu\not\ll\nu$. In the local experiments this divergence is applied to one-record observed-data margins such as $P_O$ and $(P_0^{\mathrm{wit}})_O$.
\end{definitionv}

The ordered-weight lower bound applies the Le Cam local testing framework of \citet{LeCam1986} and \citet{PolyanskiyWu2019} to \cref{obj:ass:local-weight-neighborhood}. The center has effect separation indexed by \(g\), matching the gap-local target used in \cref{obj:def:gap-stratum}. The next display repeats the same divergence notation for the local-experiment statements that follow.

The witness laws are finite Bernoulli proxy models with explicit masses, treatment probabilities, proxy features, and potential-outcome means.

\begin{definitionv}{synth_15}[Kullback--Leibler divergence]
For probability laws $\mu$ and $\nu$ on the same measurable space, the Kullback--Leibler divergence from $\nu$ to $\mu$ is
\[
D_{\mathrm{KL}}(\mu\|\nu):=\int \log\!\left(\frac{d\mu}{d\nu}\right)\,d\mu
\]
when $\mu\ll\nu$, and $D_{\mathrm{KL}}(\mu\|\nu):=+\infty$ when $\mu\not\ll\nu$. In the local experiments this divergence is applied to one-record observed-data margins, including $P_O$ and $(P_0^{\mathrm{wit}})_O$.
\end{definitionv}

\begin{assumptionv}{ass:local-quotient-neighborhood}[Local quotient neighborhood]
The Kullback--Leibler divergence from \(P_O\), the observed-data marginal, to the observed margin of \(P_0^{\mathrm{wit}}\), the two-class collision witness at perturbation zero, satisfies
\[
D_{\mathrm{KL}}\bigl(P_O\Vert (P_0^{\mathrm{wit}})_O\bigr)\le \frac{c_{\mathrm{loc}}}{n}.
\]
\end{assumptionv}

\Cref{obj:prop:two-class-witness-valid} places the collision witness inside the same uniformly conditioned class used for the upper bounds. The determinants and arm probabilities show that the testing pair is an interior proxy model under the displayed margins, while the quotient law moves from a single atom at \(\varepsilon=0\) to two nearby atoms for positive \(\varepsilon\).

For labeled weights, the path changes latent masses while preserving the proxy factorizations that determine the observable \(ZX^\top\) moments.

\begin{assumptionv}{ass:local-weight-neighborhood}[Local weight neighborhood]
The Kullback--Leibler divergence from \(P_O\), the observed-data marginal, to the observed margin of the two-class witness at perturbation \(g/2\) satisfies
\[
D_{\mathrm{KL}}\bigl(P_O\Vert (P_{g/2}^{\mathrm{wit}})_O\bigr)\le \frac{c_{\mathrm{loc}}}{n}.
\]
\end{assumptionv}

\begin{definitionv}{def:two-class-witness}[Two-class witness law \(P_{\varepsilon}^{\mathrm{wit}}\)]
For \(k=d_x=d_z=2\), the two-class witness law \(P_{\varepsilon}^{\mathrm{wit}}\) is the law specified by
\[
P(U=1)=0.4,\qquad P(U=2)=0.6,
\]
\[
P(T=1\mid U=1)=0.4,\qquad P(T=1\mid U=2)=0.6,
\]
\[
X=(1,X_2)^\top,\qquad Z=(1,Z_2)^\top,
\]
\[
B=\begin{psmallmatrix}1&1\\0.2&0.8\end{psmallmatrix},
\qquad
A_0=\begin{psmallmatrix}1&1\\0.3&0.7\end{psmallmatrix},
\qquad
A_1=\begin{psmallmatrix}1&1\\0.35&0.75\end{psmallmatrix},
\]
\[
E[Y(0)\mid U]=(0.25,0.25),
\qquad
E[Y(1)\mid U]=(0.5-\varepsilon,0.5+\varepsilon).
\]
Conditional on \(U\), the variables \(X_2\), \(Y(0)\), \(Y(1)\), and \(T\) are mutually independent Bernoulli variables with the displayed conditional means. Conditional on \(U\) and \(T\), the variable \(Z_2\) is independent of \((X_2,Y(0),Y(1))\) and has conditional mean given by \(A_T\). The observed outcome is
\[
Y=Y(T).
\]
\end{definitionv}

The family \(P_{g,h}\) supplies the local perturbation for the ordered-weight converse. At \(h=0\) it coincides with the separated witness at displacement \(g/2\), and changing \(h\) changes the ordered masses while retaining the displayed armwise proxy factorizations.

The two-class witness also gives a compact picture of the reporting geometry. In this specialization the class masses are \(0.4\) and \(0.6\), and the class-average effects are \(0.25-\varepsilon\) and \(0.25+\varepsilon\). With association radius \(\rho_{n,\alpha}\), the empirical graph in \cref{obj:def:cluster-report} connects estimated atoms at threshold \(4\rho_{n,\alpha}\); the support interval then expands each empirical component by \(\rho_{n,\alpha}\).

\[
\begin{array}{c|c|c|c}
\text{witness regime} & \text{quotient law} & \text{component rule at radius }\rho_{n,\alpha} & \text{reported mass target}\\
\hline
\varepsilon=0
& \delta_{0.25}
& \text{one component around }0.25
& 1\\
0<2\varepsilon\le 4\rho_{n,\alpha}
& 0.4\,\delta_{0.25-\varepsilon}+0.6\,\delta_{0.25+\varepsilon}
& \text{connected empirical component}
& 0.4+0.6\\
2\varepsilon>4\rho_{n,\alpha}
& 0.4\,\delta_{0.25-\varepsilon}+0.6\,\delta_{0.25+\varepsilon}
& \text{two empirical components when the center is separated}
& 0.4,\ 0.6
\end{array}
\]

The table serves as a reading guide for the existing witness. It illustrates how quotient aggregation and component reporting use the same masses while the merge threshold determines whether the report presents one aggregate interval or two separated component intervals.

\begin{definitionv}{def:local-quotient-experiment}[Local quotient experiment \(\mathcal L_n^\nu\)]
The local quotient-law KL experiment is
\[
\mathcal L_n^\nu
:=
\{P\in\mathcal M:P\text{ satisfies }\cref{obj:ass:local-quotient-neighborhood}\},
\]
where \(\mathcal M\) is the model class in \cref{obj:def:model-class}.
\end{definitionv}

\begin{definitionv}{def:local-weight-experiment}[Local weight experiment \(\mathcal L_{n,g}^p\)]
The local labeled-weight KL experiment is
\[
\mathcal L_{n,g}^p
:=
\{P\in\mathcal M(g):P\text{ satisfies }\cref{obj:ass:local-weight-neighborhood}\},
\]
where \(\mathcal M(g)\) is the gap-local stratum in \cref{obj:def:gap-stratum}.
\end{definitionv}

The two local experiments distinguish the metrics under study. \(\mathcal L_n^\nu\) is centered at a collision and measures the law-level \(W_1\) difficulty, whereas \(\mathcal L_{n,g}^p\) is centered on the gap-local stratum and measures the difficulty of recovering effect-ordered masses.

\begin{propositionv}{prop:two-class-witness-valid}[Two-class witness validity]
For every collision-witness effect displacement \(\varepsilon\) satisfying
\[
0\le \varepsilon \le \frac18,
\]
the explicit two-class Bernoulli proxy witness law \(P_{\varepsilon}^{\mathrm{wit}}\) belongs to \(\mathcal M\) with \(L=2\), \(\pi_0=0.1\), and \(\sigma_0=0.1\). Moreover,
\[
P_{\varepsilon}^{\mathrm{wit}}(T=0)=0.48,\qquad
P_{\varepsilon}^{\mathrm{wit}}(T=1)=0.52,
\]
\[
\det(B)=0.6,\qquad
\det(A_0)=\det(A_1)=0.4,
\]
\[
\det(M_0(P_{\varepsilon}^{\mathrm{wit}}))=0.06,\qquad
\det(M_1(P_{\varepsilon}^{\mathrm{wit}}))=0.4\cdot 0.6\cdot \frac{36}{169},
\]
and the quotient latent-effect law is
\[
\nu_{P_{\varepsilon}^{\mathrm{wit}}}
=
0.4\,\delta_{0.25-\varepsilon}
+
0.6\,\delta_{0.25+\varepsilon}.
\]
At \(\varepsilon=0\), the proxy feature and observable moment matrices above are nonsingular, the quotient latent-effect law is \(\delta_{0.25}\), and the two latent-class masses are distinct.
\end{propositionv}
It is worth tracing the collision through the witness once, since every object in the paper is visible on it. The latent effects of \(P_{\varepsilon}^{\mathrm{wit}}\) are read off the displayed conditional means as \(\tau_1=0.25-\varepsilon\) and \(\tau_2=0.25+\varepsilon\), carried by the class masses \(0.4\) and \(0.6\), so the effect gap is \(2\varepsilon\) and the quotient law is
\[
\nu_{P^{\mathrm{wit}}_{\varepsilon}}=0.4\,\delta_{0.25-\varepsilon}+0.6\,\delta_{0.25+\varepsilon}
\qquad(\varepsilon>0),
\qquad
\nu_{P^{\mathrm{wit}}_{0}}=\delta_{0.25}.
\]

\begin{center}
\small
\begin{tabular}{p{0.13\linewidth}p{0.20\linewidth}p{0.10\linewidth}p{0.24\linewidth}p{0.21\linewidth}}
regime & latent effects & gap & quotient law & effect-ordered masses\\
\hline
\(\varepsilon=0\) & \(0.25,\;0.25\) & \(0\) & one atom at \(0.25\) of mass \(1\) & labels carry no separation\\
\(\varepsilon>0\) & \(0.25-\varepsilon,\;0.25+\varepsilon\) & \(2\varepsilon\) & \(0.4\) at \(0.25-\varepsilon\), \(0.6\) at \(0.25+\varepsilon\) & \((0.4,0.6)\)\\
\end{tabular}
\end{center}

Transporting the mass \(0.4\) and the mass \(0.6\) each a distance \(\varepsilon\) onto the common location gives \(W_1(\nu_{P^{\mathrm{wit}}_{\varepsilon}},\nu_{P^{\mathrm{wit}}_{0}})=\varepsilon\), so the estimand moves continuously into the collision and the modulus of \cref{obj:thm:gap-free-positive-measure-modulus} applies across the whole family at fixed conditioning constants. The two regimes are distinguished by which report they support. The quotient law and its \(W_1\) confidence set of \cref{obj:def:wasserstein-confidence-set} remain meaningful at \(\varepsilon=0\), where the two classes have become one atom of mass \(1\). The effect-ordered masses \((0.4,0.6)\) are separated only while \(2\varepsilon\) is positive, which is why \cref{obj:thm:labeled-weight-upper} is stated on the gap-local strata and why its converse in \cref{obj:thm:matching-local-lower-bounds} is driven by displacement along this same family. The cluster report of \cref{obj:def:cluster-report} interpolates between the two pictures: while the estimated support points fall within the connection threshold \(4\rho_{n,\alpha}\) they are returned as one component whose aggregate mass interval covers the combined mass, and once the separation exceeds that threshold they are reported as two components whose mass intervals sharpen with the external separation, in the manner quantified by \cref{obj:thm:cluster-adaptive-report}.


\Cref{obj:thm:matching-local-lower-bounds} gives the two local converse rates on the displayed two-class specialization. For quotient laws, an \(n^{-1/2}\) displacement in the collision witness creates \(W_1\) separation of order \(n^{-1/2}\) while keeping the observed Kullback--Leibler comparison at order \(n^{-1}\). For ordered weights, the factorization-preserving path makes the observed divergence scale as \(g^2h^2\), so the largest indistinguishable mass displacement is clipped at \(\min\{1,(\sqrt n g)^{-1}\}\).

These local statements calibrate the upper bounds from the preceding sections on the same explicit two-class specialization. The quotient-law lower bound matches the root-\(n\) law rate from \cref{obj:thm:collision-uniform-root-n}; the ordered-weight lower bound matches the inverse-gap upper bound in \cref{obj:thm:labeled-weight-upper}.

\begin{theoremv}{thm:matching-local-lower-bounds}[Matching local lower bounds]
Fix \(c_{\mathrm{loc}}\in(0,1)\). There exist constants \(a,c,C\in\mathbb R\) such that
\[
0<a\le \frac18,\qquad c>0,\qquad C>0,
\]
and, for every integer \(n\ge 1\), the following statements hold for the explicit submodel \(k=d_x=d_z=2\), \(L=2\), \(\pi_0=1/10\), and \(\sigma_0=1/10\).
\begin{itemize}
\item \textbf{(Quotient-law risk.)} For every measurable quotient-law estimator \(\widetilde\nu_n\) with support radius
\[
L_\tau=\frac{4L\sqrt{d_z}}{\sigma_0},
\]
there is a law \(P\) in the local quotient-law experiment \(\mathcal L_n^\nu\) of \cref{obj:def:local-quotient-experiment} such that
\[
P=P_0^{\mathrm{wit}}
\quad\text{or}\quad
P=P_{a/\sqrt n}^{\mathrm{wit}},
\]
with \(P_\varepsilon^{\mathrm{wit}}\) as in \cref{obj:def:two-class-witness}, and
\[
\frac{c}{\sqrt n}
\le
E_{Q_P^{(n)}}\!\left[W_1(\widetilde\nu_n,\nu_P)\right].
\]

\item \textbf{(Ordered-weight risk.)} For every \(g\in(0,1/4]\) and every measurable ordered-weight estimator \(\widetilde p_n\) taking values in the simplex, there is a law \(P\) in the local labeled-weight experiment \(\mathcal L_{n,g}^{p}\) of \cref{obj:def:local-weight-experiment} such that
\[
P=P_{g,0}
\quad\text{or}\quad
P=P_{g,h(n,g)},
\qquad
h(n,g)=a\min\left\{1,\frac{1}{\sqrt n\,g}\right\},
\]
and
\[
c\min\left\{1,\frac{1}{\sqrt n\,g}\right\}
\le
E_{Q_P^{(n)}}\!\left[\|\widetilde p_n-p^{\uparrow}(P)\|_1\right].
\]

\item \textbf{(Quotient-law certificate.)} The laws \(P_0^{\mathrm{wit}}\) and \(P_{a/\sqrt n}^{\mathrm{wit}}\) belong to the corresponding local quotient-law experiments, and their observed-data marginals and quotient laws satisfy
\[
D_{\mathrm{KL}}\!\left((P_{a/\sqrt n}^{\mathrm{wit}})_O\Vert (P_0^{\mathrm{wit}})_O\right)
\le
\frac{C}{n},
\qquad
\frac{c}{\sqrt n}
\le
W_1\!\left(\nu_{P_0^{\mathrm{wit}}},\nu_{P_{a/\sqrt n}^{\mathrm{wit}}}\right).
\]

\item \textbf{(Labeled-path certificate.)} For every \(g\in(0,1/4]\), with
\[
h=h(n,g)=a\min\left\{1,\frac{1}{\sqrt n\,g}\right\},
\]
the displacement \(h\) lies in the tangent-amplitude domain, the laws \(P_{g,h}\) and \(P_{g,0}\) belong to the corresponding local labeled-weight experiments, and
\[
D_{\mathrm{KL}}\!\left((P_{g,h})_O\Vert (P_{g,0})_O\right)
\le
C g^2h^2,
\qquad
\sum_i\left|p_i^{\uparrow}(P_{g,h})-p_i^{\uparrow}(P_{g,0})\right|
=
2|h|.
\]
\end{itemize}
\end{theoremv}

\Cref{obj:prop:same-class-quotient-minimax} converts the local quotient-law comparison into a same-class minimax statement on the displayed uniformly conditioned two-class model. The proposition states that the optimal uniform \(W_1\) risk for the quotient law is of exact root-\(n\) order on this specialization.

\begin{propositionv}{prop:same-class-quotient-minimax}[Same-class quotient minimax]
For the specialization \(k=d_x=d_z=2\), \(L=2\), \(\pi_0=1/10\), and \(\sigma_0=1/10\), let \(\mathcal M\) be the uniformly conditioned proxy causal model class in \cref{obj:def:model-class}.  The derived effect-support radius is
\[
L_\tau=4L\sqrt{d_z}/\sigma_0=80\sqrt 2 .
\]
There exist real constants \(c\) and \(C\) such that
\[
0<c<C
\]
and, for every sample size \(n\ge 1\), the following two statements hold:
\begin{itemize}
\item \textbf{(Lower bound.)} For every measurable estimator \(\widetilde\nu_n\) mapping an \(n\)-sample of observed records to a represented atomic law modulo the radius \(L_\tau\), there exists a full-data probability law \(P\in\mathcal M\) such that
\[
\frac{c}{\sqrt n}
\le
E_{Q_P^{(n)}}\!\left[
  W_1\!\left(\widetilde\nu_n,\nu_P\right)
\right].
\]
\item \textbf{(Upper bound.)} There exists a measurable estimator \(\widetilde\nu_n\) mapping an \(n\)-sample of observed records to a represented atomic law modulo the radius \(L_\tau\) such that, for every full-data probability law \(P\in\mathcal M\),
\[
E_{Q_P^{(n)}}\!\left[
  W_1\!\left(\widetilde\nu_n,\nu_P\right)
\right]
\le
\frac{C}{\sqrt n}.
\]
\end{itemize}
Equivalently, on this same specialized class,
\[
\frac{c}{\sqrt n}
\le
\inf_{\widetilde\nu_n}\sup_{P\in\mathcal M}
E_{Q_P^{(n)}}\!\left[
  W_1\!\left(\widetilde\nu_n,\nu_P\right)
\right]
\le
\frac{C}{\sqrt n}
\]
for every \(n\ge 1\), where the infimum ranges over measurable estimators with output radius \(L_\tau\).
\end{propositionv}

\Cref{obj:prop:same-class-labeled-minimax} gives the corresponding same-class minimax statement for ordered weights on the same two-class specialization. The rate is governed by the product \(\sqrt n g\): when the effect gap is large relative to sampling noise, ordered masses are estimable at the inverse-gap scale, and the bound remains clipped by the diameter of the simplex-valued loss.

The final comparison records how the same witnesses transfer to classes formulated in the terminology of \citet{VirkMazaheriWu2026}. The following definitions separate the published structural gates, finite-regularity condition, and separated recovery regime used for that transfer.

\begin{propositionv}{prop:same-class-labeled-minimax}[Same-class labeled minimax]
There exist real constants \(c\) and \(C\) such that \(0<c<C\) and, for every sample size \(n\ge 1\) and every gap scale \(g\) with \(0<g\le 1/4\), the gap stratum \(\mathcal M(g)\) from \cref{obj:def:gap-stratum} satisfies
\[
c\min\left\{1,\frac{1}{\sqrt n\,g}\right\}
\le
\inf_{\widetilde p_n}
\sup_{P\in\mathcal M(g)}
E_{Q_P^{(n)}}\!\left[
\left\|\widetilde p_n-p^{\uparrow}(P)\right\|_1
\right]
\le
C\min\left\{1,\frac{1}{\sqrt n\,g}\right\}.
\]
Here the infimum is over all measurable simplex-valued labeled weight estimators based on \(n\) observed records in the explicit submodel \(k=d_x=d_z=2\), \(L=2\), \(\pi_0=1/10\), and \(\sigma_0=1/10\), and \(p^{\uparrow}(P)\) is the effect-ordered latent-mass vector with the barycentric collision fallback.
\end{propositionv}

\begin{definitionv}{synth_31}[Calibrated path displacement]
For \(a,g\in\mathbb R\) and \(n\in\mathbb N\), define
\[
\operatorname{calibratedDisplacement}(a,n,g)
=
a\,\min\{1,(\sqrt n\,g)^{-1}\}.
\]
\end{definitionv}

\begin{definitionv}{synth_12}[Published VMW finite-regularity condition]
A law \(P\) satisfies the published VMW finite-regularity condition when \(0<k\), \(k\le d_x\), \(k\le d_z\), and there exist constants \(L_X,L_{ZX},L_{YZX},\pi,\sigma>0\) such that \(\lVert X\rVert_2\le L_X\), \(\lVert ZX^\top\rVert_{\mathrm{op}}\le L_{ZX}\), and \(\lVert YZX^\top\rVert_{\mathrm{op}}\le L_{YZX}\) almost surely, \(P(T=t)\ge\pi\) for \(t\in\{0,1\}\), and there is a top-\(k\) right singular basis \(V\) for the stacked proxy moment \([M_0(P);M_1(P)]\) satisfying \(\sigma_k([M_0(P);M_1(P)])\ge\sigma\) and \(\sigma_k(M_t(P)V)\ge\sigma\) for each \(t\in\{0,1\}\).
\end{definitionv}

\begin{definitionv}{synth_11}[Published VMW structural model]
Fix \(k,d_x,d_z\) with \(0<k\), \(k\le d_x\), and \(k\le d_z\). A one-unit full-data probability law \(P\) belongs to the published VMW structural model when it satisfies the published qualitative proxy model conditions: reference-proxy separation \(Z\perp (X,Y)\mid(T,U)\), target-proxy separation \(X\perp (Y,T)\mid U\), consistency \(Y=Y(T)\) almost surely, armwise latent ignorability \(Y(t)\perp T\mid U\) for \(t\in\{0,1\}\), full column rank of \(A_0\), \(A_1\), and \(B\), and strict latent-arm positivity \(P(U=u,T=t)>0\) for every \(u\in\{1,\ldots,k\}\) and \(t\in\{0,1\}\).
\end{definitionv}

\Cref{obj:def:published-vmw-regimes}, \cref{obj:synth_12}, \cref{obj:synth_11}, and \cref{obj:synth_13} provide the interface for comparing the present lower bounds with the published VMW recovery framework. The cited model conditions and separated recovery regime are used as external publication-level gates, while the converse below is stated in terms of exact containment of the two displayed witness pairs.

\begin{definitionv}{synth_13}[Separated recovery regime]
A full-data law \(P\) lies in the separated recovery regime when it satisfies:
\begin{itemize}
\item \textbf{(Structural model.)} The published VMW structural model conditions.
\item \textbf{(Anchor normalization.)} The first target-proxy coordinate satisfies \(X_1=1\) almost surely.
\item \textbf{(Finite regularity.)} The published VMW finite-regularity condition.
\item \textbf{(Spectral separation.)} The latent-class treatment effects satisfy \(\tau_u\ne\tau_v\) whenever \(u\ne v\).
\end{itemize}
In this regime, the separated-recovery estimator analyzed by \citet{VirkMazaheriWu2026} is evaluated under its stated sample-size and radius conditions and returns simultaneous high-probability bounds for the treatment effects, the anchor-normalized feature matrix, and the simplex-projected mixture weights.
\end{definitionv}

The calibrated displacement in \cref{obj:synth_31} is the path amplitude used in the labeled-weight two-point comparison. It is the same clipped scale that appears in the ordered-weight upper and lower bounds.

\begin{theoremv}{thm:published-vmw-converse-transfer}[Published converse transfer]*
Fix \(k=d_x=d_z=2\), \(L=2\), and \(\sigma_0=1/10\).  There exist constants \(a,c\in\mathbb R\) such that
\[
0<a\le \frac18,
\qquad
0<c .
\]
For every \(n\ge 1\), the following hold.
\begin{itemize}
\item \textbf{(Quotient witness pair.)} The two laws \(P_0^{\mathrm{wit}}\) and \(P_{a/\sqrt n}^{\mathrm{wit}}\) from \cref{obj:def:two-class-witness} belong to the published VMW model and satisfy VMW Assumption 4, while \(P_0^{\mathrm{wit}}\) lies outside the published separated recovery regime.
\item \textbf{(Quotient lower bound.)} For every class \(\mathcal V_n^\nu\) of full-data probability laws such that every \(P\in\mathcal V_n^\nu\) belongs to the published VMW model and satisfies VMW Assumption 4, if \(\mathcal V_n^\nu\) contains laws equal to \(P_0^{\mathrm{wit}}\) and \(P_{a/\sqrt n}^{\mathrm{wit}}\), then every estimator \(\widetilde\nu_n\) of the quotient effect law with support radius \(4\cdot 2\sqrt2/(1/10)\) has worst-case risk at least \(c/\sqrt n\):
\[
\sup_{P\in\mathcal V_n^\nu}
E_{Q_P^{(n)}}\!\left[
W_1\!\left(\widetilde\nu_n,\nu_P\right)
\right]
\ge \frac{c}{\sqrt n}.
\]
\item \textbf{(Separated labeled path.)} For every \(g\) in the gap-scale domain, define
\[
h(n,g):=\operatorname{calibratedDisplacement}(a,n,g).
\]
Then \(h(n,g)\) is in the tangent-amplitude domain, and the endpoint laws \(P_{g,0}\) and \(P_{g,h(n,g)}\) both belong to the published VMW model, satisfy VMW Assumption 4, and have spectrally separated latent effects.
\item \textbf{(Labeled-weight lower bound.)} For every \(g\) in the gap-scale domain and every class \(\mathcal V_{n,g}^p\) of full-data probability laws such that every \(P\in\mathcal V_{n,g}^p\) belongs to the published VMW model and satisfies VMW Assumption 4, if \(\mathcal V_{n,g}^p\) contains laws equal to \(P_{g,0}\) and \(P_{g,h(n,g)}\), then every estimator \(\widetilde p_n\) of the ordered latent weights satisfies
\[
\sup_{P\in\mathcal V_{n,g}^p}
E_{Q_P^{(n)}}\!\left[
\left\|\widetilde p_n-p^{\uparrow}(P)\right\|_1
\right]
\ge
c\min\{1,(\sqrt n\,g)^{-1}\}.
\]
\end{itemize}
\end{theoremv}
% CAUSALSMITH-CITED-SCOPE-BEGIN thm:published-vmw-converse-transfer
\verificationfootnotetext{\textbf{Formalization scope.} This result uses the published conclusion \citep{VirkMazaheriWu2026} (Assumptions 1--2 and Section 4.2 of arXiv:2607.10926v1) and \citep{VirkMazaheriWu2026} (Assumption 3, lines 261--263; Assumption 4, lines 321--337; and Theorem 7.2, lines 350--381 of arXiv:2607.10926v1). The cited source proof is not formalized here: Lean verifies its use and all remaining steps. If that cited conclusion is false, the portions of this result that depend on it are not certified.}
% CAUSALSMITH-CITED-SCOPE-END thm:published-vmw-converse-transfer

\Cref{obj:thm:published-vmw-converse-transfer} states the converse comparison with \citet{VirkMazaheriWu2026}. Any comparator class satisfying the cited VMW model and finite-regularity gates inherits the displayed lower bound once it contains the two quotient witness laws or the two labeled-path endpoint laws. The transfer is therefore tied to concrete witness containment: the quotient comparison uses the collision pair, while the labeled-weight comparison uses separated endpoint laws at the calibrated displacement.
